\chapter{Data Visualization \& Course Summary}

Data Visualization turns raw numerical data into intuitive charts, graphs, and interactive dashboards. In this final chapter, we cover \textbf{Matplotlib} (the fundamental plotting library in Python) and review the core programming concepts mastered throughout the course.

\section{Introduction to Matplotlib}

\lstinline|matplotlib.pyplot| provides a MATLAB-like plotting framework for building line plots, scatter plots, bar charts, histograms, and heatmaps.

\begin{definitionbox}{Matplotlib Figure and Axes}
A \textbf{Figure} represents the entire canvas window. An \textbf{Axes} represents an individual plot or subplot inside the figure.
\begin{lstlisting}
import matplotlib.pyplot as plt
import numpy as np

x = np.linspace(0, 10, 100)
y = np.sin(x)

plt.figure(figsize=(8, 4))
plt.plot(x, y, label="Sine Wave", color="blue", linestyle="--")
plt.title("Trigonometric Sine Function")
plt.xlabel("X values")
plt.ylabel("sin(x)")
plt.legend()
plt.grid(True)
plt.show()
\end{lstlisting}
\end{definitionbox}

\begin{videocard}{IITM BS Lecture: Introduction to Matplotlib}{https://www.youtube.com/watch?v=hk6NsCfvftY}
Official lecture demonstrating plot creation, line charts, scatter plots, titles, labels, and legends.
\end{videocard}

\section{Common Chart Types for Data Science}

\begin{theorembox}{Plot Types and Use Cases}
\begin{itemize}
    \item \textbf{Line Plot (\lstinline|plt.plot|):} Trends over continuous time or ordered variables.
    \item \textbf{Scatter Plot (\lstinline|plt.scatter|):} Relationship and correlation between two numerical variables.
    \item \textbf{Bar Chart (\lstinline|plt.bar|):} Comparing numerical measurements across discrete categories.
    \item \textbf{Histogram (\lstinline|plt.hist|):} Frequency distribution of a single continuous variable.
\end{itemize}
\end{theorembox}

\begin{examplebox}{Scatter Plot and Histogram Example}
\begin{lstlisting}
# Scatter Plot for Feature Association
heights = [160, 165, 170, 175, 180]
weights = [55, 62, 68, 74, 82]
plt.scatter(heights, weights, color='purple')
plt.title("Height vs Weight Scatter Plot")

# Histogram for Distribution Analysis
data = np.random.normal(loc=0, scale=1, size=1000)
plt.hist(data, bins=30, color='teal', edgecolor='black')
plt.title("Normal Distribution Histogram")
\end{lstlisting}
\end{examplebox}

\index{Matplotlib}\index{Scatter Plot}\index{Histogram}\index{EDA}

\textbf{Data Science Relevance:} Data visualization is critical during Exploratory Data Analysis (EDA) to detect outliers, assess skewness, inspect feature correlations, and present findings to business stakeholders.


\section{Course Summary and Next Steps}

Congratulations! You have completed the foundational Python programming curriculum for Data Science.

\begin{formulabox}{Core Curriculum Overview}
\begin{enumerate}
    \item \textbf{Weeks 1–3:} Variables, Types, Control Flow, and Loops.
    \item \textbf{Weeks 4–5:} Data Structures (Lists, Tuples, Dicts, Sets) and Matrix Operations.
    \item \textbf{Weeks 6–7:} Modular Design with Functions, Scopes, Comprehensions, and Generators.
    \item \textbf{Weeks 8–9:} File Processing, Robust Exception Handling, Recursion, and Binary Search.
    \item \textbf{Weeks 10–12:} Object-Oriented Architecture, Scientific Computing (NumPy/Pandas), and Visualization (Matplotlib).
\end{enumerate}
\end{formulabox}

\begin{videocard}{IITM BS Lecture: Course Summary and Next Steps}{https://www.youtube.com/watch?v=XYDiGqzVKyo}
Official course summary reviewing core concepts and outlining advanced topics in Data Science.
\end{videocard}

\newpage
\section{Practice Exercises}
\textit{Detailed step-by-step solutions are available in the Appendix.}

\begin{enumerate}
    \item \textbf{Basic Line Plot:} Write a Matplotlib script to plot $y = x^2$ for $x \in [-5, 5]$ with $100$ points. Add title \lstinline|"Quadratic Function"|, x-label \lstinline|"x"|, y-label \lstinline|"f(x) = x^2"|, and display a grid.

    \item \textbf{Comparing Distributions (Subplots):} Write a Matplotlib script using \lstinline|plt.subplots(1, 2)| to display two side-by-side plots:
    \begin{enumerate}
        \item Left Subplot: A bar chart of top 3 programming languages \lstinline|['Python', 'SQL', 'R']| with popularity scores \lstinline|[95, 80, 60]|.
        \item Right Subplot: A histogram of $500$ random points drawn from a uniform distribution between $0$ and $100$.
    \end{enumerate}

    \item \textbf{Customizing Scatter Plots:} Create a scatter plot of $20$ random $(x, y)$ points. Set point color to green, marker style to \lstinline|'^'| (triangle), and marker size to $50$.

    \item \textbf{Comprehensive Review Question (Data Pipeline):} Write a Python function \lstinline|analyze_dataset(data)| that takes a list of raw numeric measurements:
    \begin{enumerate}
        \item Cleans the list by removing negative values using list comprehension.
        \item Calculates the mean and standard deviation of the valid numbers using NumPy.
        \item Returns a dictionary \lstinline|{"count": N, "mean": M, "std": S}|.
    \end{enumerate}

    \item \textbf{Data Science Application (Plotting a Loss Curve):} In deep learning, model loss decreases over training epochs. Given \lstinline|epochs = range(1, 11)| and \lstinline|loss = [2.5, 1.8, 1.2, 0.8, 0.5, 0.35, 0.25, 0.18, 0.15, 0.12]|, write a Matplotlib script to plot the training loss curve with a red line and square markers.
\end{enumerate}